\chapter[Galois Theory]{Galois Theory, Limits, and Profinite Groups}

Galois theory replaces the study of intermediate fields by the study of the
symmetries of the larger field.  In the finite case this is an exact
dictionary.  In the infinite case the same dictionary survives once the
symmetry group remembers agreement at every finite stage through a topology.

\section{Galois Groups and Fixed Fields}

\begin{definition}
For an extension \(E/F\), its \emph{Galois group} is
\[
\operatorname{Gal}(E/F)=\operatorname{Aut}_F(E).
\]
If \(H\leq\operatorname{Gal}(E/F)\), its \emph{fixed field} is
\[
E^H=\{x\in E:\sigma(x)=x\text{ for every }\sigma\in H\}.
\]
For an intermediate field \(F\subseteq K\subseteq E\), write
\[
\operatorname{Gal}(E/K)=\{\sigma\in\operatorname{Gal}(E/F):
\sigma|_K=1_K\}.
\]
\end{definition}

The fixed field records the elements on which the chosen symmetries are
invisible; it will be the inverse operation to taking an automorphism group.

\begin{lemma}\label{lem:fixed-field-intermediate}
The set \(E^H\) is an intermediate field.  Furthermore,
\[
K_1\subseteq K_2\Longrightarrow
\operatorname{Gal}(E/K_2)\leq\operatorname{Gal}(E/K_1),
\qquad
H_1\leq H_2\Longrightarrow E^{H_2}\subseteq E^{H_1}.
\]
\end{lemma}
\begin{proof}
Every element of \(H\) fixes \(0\) and \(1\).  If \(x,y\in E^H\), then each
\(\sigma\in H\) fixes \(x-y\), \(xy\), and \(x^{-1}\) when \(x\ne0\).
Thus \(E^H\) is a subfield containing \(F\).  The two inclusions say,
respectively, that fixing more elements imposes more conditions and that
asking more automorphisms to fix an element imposes more conditions.
\end{proof}

\begin{example}
Complex conjugation has fixed field \(\mathbb R\) in \(\mathbb C\).  If \(d\)
is squarefree, the nonidentity automorphism of \(\mathbb Q(\sqrt d)\) sends
\(\sqrt d\) to \(-\sqrt d\) and fixes precisely \(\mathbb Q\).  In
\(\mathbb Q(\zeta_5)\), the automorphisms are
\(\zeta_5\mapsto\zeta_5^a\), \(a\in(\mathbb Z/5\mathbb Z)^\times\); its
unique order-two subgroup fixes \(\mathbb Q(\sqrt5)\).
\end{example}

The following linear-algebra fact is the key input in Artin's theorem.
Distinct field homomorphisms cannot satisfy an accidental linear relation as
functions.

\begin{lemma}[Artin independence lemma]\label{lem:artin-independence}
Let \(\sigma_1,\ldots,\sigma_n:E\to L\) be distinct field homomorphisms.
Then they are linearly independent over \(L\) as functions \(E\to L\).
\end{lemma}
\begin{proof}
Suppose that a nontrivial relation
\[
a_1\sigma_1+\cdots+a_n\sigma_n=0
\]
exists, and choose one with the fewest nonzero coefficients.  Omit zero
terms, so every \(a_i\ne0\) and \(n\ge2\).  Choose \(x\in E\) with
\(\sigma_1(x)\ne\sigma_n(x)\).  Evaluating at \(xy\), then subtracting
\(\sigma_n(x)\) times the original relation evaluated at \(y\), yields
\[
\sum_{i=1}^{n-1}a_i\bigl(\sigma_i(x)-\sigma_n(x)\bigr)\sigma_i(y)=0
\quad\text{for all }y\in E.
\]
The coefficient of \(\sigma_n\) has vanished while that of \(\sigma_1\) is
nonzero.  This is a shorter nontrivial relation, a contradiction.
\end{proof}

The proof of Artin's theorem has a useful two-sided shape.  Independence of
the automorphisms first produces enough linearly independent elements to give
the lower bound \(|G|\leq[E:E^G]\).  Orbit polynomials then show that every
finitely generated subfield has degree at most \(|G|\), producing the reverse
bound.  The final embedding count identifies the automorphisms.  This is the
bridge from a concrete group action to the field--subgroup dictionary.

\begin{theorem}[Artin fixed-field theorem]\label{thm:artin-fixed-field}
If \(G\) is a finite group of automorphisms of a field \(E\), then
\[
[E:E^G]=|G|,
\qquad
\operatorname{Gal}(E/E^G)=G.
\]
\end{theorem}
\begin{proof}
Put \(F=E^G\) and \(n=|G|\).  Artin independence implies that there are
\(x_1,\ldots,x_n\in E\) for which the evaluation matrix
\[
\bigl(\sigma_i(x_j)\bigr)_{1\leq i,j\leq n}
\]
is invertible.  Indeed, consider the evaluation vectors
\[
v(x)=(\sigma_1(x),\ldots,\sigma_n(x))\in E^n.
\]
If these did not span \(E^n\), a nonzero linear functional
\((a_1,\ldots,a_n)\) would vanish on every \(v(x)\), giving
\[
a_1\sigma_1+\cdots+a_n\sigma_n=0,
\]
contrary to Artin independence.  Choose \(n\) evaluation vectors forming a
basis, and use them as the columns of the displayed matrix.  The \(x_j\) are
therefore \(F\)-linearly independent, so \(n\leq[E:F]\).

For the reverse inequality, let \(\alpha\in E\) and form
\[
P_\alpha(X)=\prod_{\sigma\in G}(X-\sigma(\alpha)).
\]
Every \(\tau\in G\) permutes these factors, so the coefficients belong to
\(F\).  The minimal polynomial of \(\alpha\) over \(F\) divides
\(P_\alpha\), and is separable of degree at most \(n\).  Given finitely many
elements of \(E\), the field they generate over \(F\) is finite separable;
the primitive element theorem makes it \(F(\alpha)\), hence of degree at most
\(n\).  Therefore \(E\) has no \(n+1\) linearly independent elements over
\(F\), and \([E:F]\leq n\).

Finally \(G\leq\operatorname{Gal}(E/F)\), while the embedding-count
inequality (Theorem~\ref{thm:embedding-count}) gives
\[
|\operatorname{Gal}(E/F)|\leq[E:F]=|G|.
\]
Therefore the inclusion is equality.
\end{proof}

\begin{corollary}\label{cor:finite-automorphism-consequences}
If \(E/F\) is finite Galois, then
\[
E^{\operatorname{Gal}(E/F)}=F.
\]
Conversely, if \(G\) is a finite subgroup of \(\operatorname{Aut}(E)\), then
\(E/E^G\) is finite Galois with Galois group \(G\).
\end{corollary}
\begin{proof}
The first assertion follows by combining Artin's theorem with the finite
Galois criterion (Proposition~\ref{prop:finite-galois-criteria}).  The second
follows from the same two results.
\end{proof}

\section{The Fundamental Theorem of Finite Galois Theory}

\begin{theorem}[Finite Galois correspondence]\label{thm:finite-galois}
Let \(E/F\) be finite Galois and put \(G=\operatorname{Gal}(E/F)\).  The maps
\[
K\longmapsto\operatorname{Gal}(E/K),
\qquad
H\longmapsto E^H
\]
are inverse inclusion-reversing bijections between intermediate fields
\(F\subseteq K\subseteq E\) and subgroups \(H\leq G\).  Moreover,
\[
[E:K]=|\operatorname{Gal}(E/K)|,
\qquad
[K:F]=[G:\operatorname{Gal}(E/K)].
\]
\end{theorem}
\begin{proof}
For an intermediate \(K\), separability and normality descend from \(E/F\)
to \(E/K\): a minimal polynomial over \(K\) divides the corresponding
polynomial over \(F\), so it is separable and splits in \(E\).  Thus \(E/K\)
is finite Galois, and Corollary~\ref{cor:finite-automorphism-consequences}
gives
\[
E^{\operatorname{Gal}(E/K)}=K.
\]
For \(H\leq G\), Artin's theorem applied to the finite group \(H\) gives
\[
\operatorname{Gal}(E/E^H)=H.
\]
These identities prove the bijection.  The first degree formula is Artin's
theorem applied to \(E/K\), and the Tower Law gives
\[
[K:F]=\frac{[E:F]}{[E:K]}
=\frac{|G|}{|\operatorname{Gal}(E/K)|}.
\]
\end{proof}

The reversal is the central visual idea: a larger field has fewer
automorphisms that fix it.
\[
\begin{tikzcd}[column sep=large,row sep=small]
E \arrow[d,phantom,"\longleftrightarrow" description] & \{1\} \\
K \arrow[u] \arrow[d,phantom,"\longleftrightarrow" description] & H \arrow[u] \\
F \arrow[u] & G \arrow[u]
\end{tikzcd}
\]

\begin{proposition}[Normality refinement]\label{prop:finite-normality-refinement}
For an intermediate field \(K\),
\[
K/F\text{ is Galois}
\quad\Longleftrightarrow\quad
\operatorname{Gal}(E/K)\trianglelefteq G.
\]
When these conditions hold, restriction induces
\[
G/\operatorname{Gal}(E/K)\cong\operatorname{Gal}(K/F).
\]
\end{proposition}
\begin{proof}
Write \(H=\operatorname{Gal}(E/K)\).  First,
\[
K/F\text{ is normal}\quad\Longleftrightarrow\quad
\sigma(K)=K\text{ for every }\sigma\in G. \tag{12.1}
\]
Normality makes every conjugate of an element of \(K\) remain in \(K\).
Conversely, extend any \(F\)-embedding of \(K\) to \(E\) by
Theorem~\ref{thm:embedding-extension}; normality of \(E/F\) makes the
extension an automorphism, and stability then proves normality of \(K/F\).

For \(\sigma\in G\),
\[
\sigma H\sigma^{-1}=\operatorname{Gal}(E/\sigma(K)).
\]
The finite correspondence and (12.1) show that \(H\) is normal exactly when
\(K/F\) is normal.  Every intermediate extension is separable, so this is
equivalent to Galoisness.

When \(K\) is stable, restriction
\[
\operatorname{res}:G\longrightarrow\operatorname{Gal}(K/F)
\]
is well-defined and has kernel \(H\).  An \(F\)-automorphism of \(K\) extends
to an \(F\)-embedding of \(E\), and normality of \(E/F\) makes the extension
an automorphism.  Thus restriction is onto, and the First Isomorphism Theorem
gives the quotient isomorphism.
\end{proof}

\begin{example}[The full lattice for a cubic splitting field]
Put
\[
r=\sqrt[3]2,
\qquad
E=\mathbb Q(r,\zeta_3).
\]
The roots of \(x^3-2\) are \(r,\zeta_3r,\zeta_3^2r\), and
\(\operatorname{Gal}(E/\mathbb Q)\cong S_3\) acts by permuting them.  Let
\[
K_1=\mathbb Q(r),\qquad
K_2=\mathbb Q(\zeta_3r),\qquad
K_3=\mathbb Q(\zeta_3^2r).
\]
The intermediate fields have the following lattice; each displayed
intermediate field is contained in \(E\), and the arrows are inclusions.
\[
\begin{tikzcd}[column sep=large,row sep=large]
&& E &&\\
K_1 \arrow[urr] & K_2 \arrow[ur] & K_3 \arrow[u] &
\mathbb Q(\sqrt{-3}) \arrow[ul] \\
&& \mathbb Q \arrow[ull] \arrow[ul] \arrow[u] \arrow[ur] &&
\end{tikzcd}
\]
The order-two subgroup stabilizing the indicated root fixes the corresponding
\(K_i\), so the three order-two subgroups correspond to the three cubic
subfields.  The alternating subgroup \(A_3\) fixes
\(\mathbb Q(\sqrt{-3})=\mathbb Q(\zeta_3)\), and the trivial and full
subgroups correspond to \(E\) and \(\mathbb Q\), respectively:
\[
\begin{array}{c|c|c}
\text{field (degree)} & \text{subgroup} & \text{order}\\ \hline
E\ (6) & \{1\} & 1\\
K_1,\ K_2,\ K_3\ (3) & \text{three root stabilizers} & 2\\
\mathbb Q(\sqrt{-3})\ (2) & A_3 & 3\\
\mathbb Q\ (1) & S_3 & 6
\end{array}
\]
The three order-two subgroups are not normal.  Accordingly the three cubic
fields are not Galois over \(\mathbb Q\), whereas
\(\mathbb Q(\sqrt{-3})/\mathbb Q\) is Galois.  This is a useful contrast
between the field lattice itself and the additional normality information
carried by the group lattice.
\end{example}

\begin{example}[A biquadratic contrast]
For
\[
E=\mathbb Q(\sqrt2,\sqrt3),
\]
every \(\mathbb Q\)-automorphism independently changes the signs of
\(\sqrt2\) and \(\sqrt3\).  Hence
\[
\operatorname{Gal}(E/\mathbb Q)\cong C_2\oplus C_2.
\]
Its three order-two subgroups fix, respectively,
\[
\mathbb Q(\sqrt2),\qquad
\mathbb Q(\sqrt3),\qquad
\mathbb Q(\sqrt6).
\]
Unlike the \(S_3\) example, every subgroup is normal because the group is
abelian.  Thus all three quadratic intermediate fields are Galois over
\(\mathbb Q\).  The two examples show that the same-looking degree pattern
can have very different normality behavior.
\end{example}

\begin{example}[A cyclic correspondence]
\(\operatorname{Gal}(\mathbb Q(\zeta_5)/\mathbb Q)\cong C_4\).  Its unique
subgroup of order two fixes \(\mathbb Q(\sqrt5)\).  All subgroups are normal,
so every intermediate field is Galois over \(\mathbb Q\).
\end{example}

\subsection*{Optional Application: A Galois-Theoretic Proof of the Fundamental Theorem of Algebra}

Chapter~8 used the Fundamental Theorem of Algebra to obtain an eigenvalue of a
complex operator.  We now discharge that dependence using the field and group
theory developed so far.  Apart from the standard elementary properties of
the real numbers, including square roots of nonnegative reals, the only
substantive real-analysis input is the Intermediate Value Theorem.  The
needed facts are developed in \emph{Lecture Notes on Mathematical Analysis}
\cite{MA}: polynomials are continuous (\S5.2), and the Intermediate Value
Theorem is Theorem~5.14 (\S5.4).  Hence every real polynomial of odd degree
has a real root.  Indeed, if
\[
p(x)=a_{2m+1}x^{2m+1}+\cdots,
\qquad
a_{2m+1}\ne0,
\]
then, for sufficiently large \(R>0\), leading-term behavior makes \(p(R)\)
and \(p(-R)\) have opposite signs.  Continuity and the Intermediate Value
Theorem therefore give a point \(c\in(-R,R)\) with \(p(c)=0\).

\begin{applicationtheorem}[Fundamental Theorem of Algebra]
Every nonconstant polynomial in \(\mathbb C[x]\) has a root in \(\mathbb C\).
Equivalently, \(\mathbb C\) is algebraically closed.
\end{applicationtheorem}
\begin{proof}
We first show that every finite extension of \(\mathbb R\) has degree a power
of \(2\).  Let \(M/\mathbb R\) be finite.  Characteristic zero makes
\(M/\mathbb R\) separable, so the
\hyperref[thm:finite-galois-closure]{finite Galois closure theorem} supplies
a finite Galois extension \(L/\mathbb R\) containing \(M\).  Put
\(G=\operatorname{Gal}(L/\mathbb R)\).
If an odd prime divided \(|G|\), let \(P\) be a Sylow \(2\)-subgroup.  Then
\[
[G:P]
\]
is odd and greater than one.  By the finite Galois correspondence,
\(L^P/\mathbb R\) has this odd degree.  The primitive element theorem writes
\(L^P=\mathbb R(\beta)\), so the minimal polynomial of \(\beta\) has odd
degree.  It has a root in \(\mathbb R\) by the stated external fact, which
contradicts irreducibility unless the degree is one.  Thus \(G\) is a
\(2\)-group.  Since
\[
[M:\mathbb R]\mid[L:\mathbb R]=|G|,
\]
the claimed power-of-two conclusion follows.

Fix an algebraic closure \(\overline{\mathbb R}\) and an element
\(i\in\overline{\mathbb R}\) with \(i^2=-1\), so that
\(\mathbb C=\mathbb R(i)\).  Every nontrivial quadratic extension of
\(\mathbb R\) is this field.  Indeed, if \(\mathbb R(\alpha)/\mathbb R\) is
quadratic, the minimal polynomial of \(\alpha\) is an irreducible quadratic
with negative discriminant.  Completing the square therefore produces an
element of the extension whose square is \(-1\): after division by the
leading coefficient, a polynomial \(x^2+bx+c\) gives
\[
\left(\alpha+\frac b2\right)^2=\frac{b^2-4c}{4}<0,
\]
and rescaling by the positive real square root of its negative gives a square
root of \(-1\).  It is \(i\) or \(-i\), so the extension is
\(\mathbb R(i)=\mathbb C\).

We next show directly that every element of \(\mathbb C\) has a square root.
For \(z=a+bi\), put
\[
r=\sqrt{a^2+b^2}.
\]
If \(r+a>0\), set
\[
x=\sqrt{\frac{r+a}{2}},
\qquad
y=\frac{b}{2x}.
\]
Then \(x^2-y^2=a\) and \(2xy=b\), so \((x+iy)^2=z\).  If \(r+a=0\), then
\(b=0\) and \(a\leq0\), and \((i\sqrt{-a})^2=z\).  Consequently the
quadratic formula shows that every quadratic polynomial over \(\mathbb C\)
has a root.  Hence \(\mathbb C\) has no nontrivial quadratic extension.

We claim that no finite Galois extension \(L/\mathbb R\) has degree greater
than \(2\).  Otherwise its group \(G\) is a \(2\)-group of order greater than
two.  Every nontrivial finite \(2\)-group has a normal subgroup of index two:
for order two this is clear, and in general the center theorem and Cauchy's
theorem supply a central subgroup of order two; induction in the quotient
then gives the asserted subgroup.  Choose a normal subgroup \(H\trianglelefteq
G\) of index two.  The finite correspondence makes
\[
L^H/\mathbb R
\]
a quadratic extension, hence \(L^H=\mathbb C\).  Now \(H=
\operatorname{Gal}(L/\mathbb C)\) is a nontrivial \(2\)-group, so it has a
subgroup \(J\) of index two.  The finite correspondence for \(L/\mathbb C\)
then makes
\[
L^J/\mathbb C
\]
a quadratic extension, a contradiction.  Therefore \([L:\mathbb R]\leq2\).

Finally, let \(\alpha\in\overline{\mathbb R}\).  The finite extension
\(\mathbb R(\alpha)/\mathbb R\) is separable, so the
\hyperref[thm:finite-galois-closure]{finite Galois closure theorem} gives a
finite Galois extension containing it, whose degree is at most two by the
preceding argument.  The Tower Law therefore
gives \([\mathbb R(\alpha):\mathbb R]\leq2\), so it is either \(\mathbb R\)
or, in the quadratic case, \(\mathbb C\).  Thus every \(\alpha\) belongs to \(\mathbb C\), and
\(\overline{\mathbb R}=\mathbb C\).  Hence \(\mathbb C\) is algebraically
closed.
\end{proof}

\begin{unremark}
Chapter~8 met the Fundamental Theorem of Algebra through complex linear
algebra, while its classical proof belongs to complex analysis.  The present
proof instead reveals field extensions, finite \(2\)-groups, and the Galois
correspondence behind the same theorem.
\end{unremark}

\section{Radicals and Classical Constructions}

\begin{definition}
A \emph{radical extension} is an extension contained in a tower
\[
F=K_0\subseteq K_1\subseteq\cdots\subseteq K_r,
\]
where \(K_i=K_{i-1}(\alpha_i)\) and
\(\alpha_i^{n_i}\in K_{i-1}\) for some \(n_i>0\).  A polynomial is
\emph{solvable by radicals} if its splitting field is contained in a radical
extension of its coefficient field.
\end{definition}

Roots of unity govern the conjugates of an \(n\)-th root, which explains the
following honest part of the radicals--groups connection.

\begin{proposition}[A root-of-unity radical direction]\label{prop:radical-solvable}
Let
\[
F=K_0\subseteq K_1\subseteq\cdots\subseteq K_r
\]
be a tower such that every \(K_i/F\) is finite Galois and
\[
K_i=K_{i-1}(\alpha_i),
\qquad
\alpha_i^{n_i}\in K_{i-1}.
\]
Assume that \(K_{i-1}\) contains all \(n_i\)-th roots of unity.  Then
\[
\operatorname{Gal}(K_r/F)
\]
is solvable.
\end{proposition}
\begin{proof}
If \(\alpha_i=0\), then \(K_i=K_{i-1}\), so the relevant Galois group is
trivial.  Otherwise define
\[
\chi_i:\operatorname{Gal}(K_i/K_{i-1})\longrightarrow\mu_{n_i},
\qquad
\chi_i(\sigma)=\frac{\sigma(\alpha_i)}{\alpha_i}.
\]
The equation \(\alpha_i^{n_i}\in K_{i-1}\) shows that \(\chi_i(\sigma)\)
is an \(n_i\)-th root of unity.  Since \(\mu_{n_i}\subseteq K_{i-1}\), every
\(\sigma\) fixes it pointwise; consequently
\[
\chi_i(\sigma\tau)
=\frac{\sigma\tau(\alpha_i)}{\alpha_i}
=\frac{\sigma\bigl(\chi_i(\tau)\alpha_i\bigr)}{\alpha_i}
=\chi_i(\tau)\chi_i(\sigma).
\]
The target is abelian, so this is the homomorphism law.  Moreover,
\(\chi_i(\sigma)=1\) means that \(\sigma\) fixes \(\alpha_i\); since
\(K_i=K_{i-1}(\alpha_i)\), it is then the identity.  Thus \(\chi_i\) is
injective.  Therefore \(\operatorname{Gal}(K_i/K_{i-1})\) is isomorphic to
a subgroup of the abelian group \(\mu_{n_i}\), and hence is abelian.

Put
\[
G_i=\operatorname{Gal}(K_r/K_i).
\]
The hypothesis that \(K_i/F\) is Galois makes \(G_i\) normal in
\(G_0=\operatorname{Gal}(K_r/F)\).  The finite normality refinement
(Proposition~\ref{prop:finite-normality-refinement}) gives
\[
G_{i-1}/G_i\cong\operatorname{Gal}(K_i/K_{i-1}).
\]
Consequently
\[
G_0\trianglerighteq G_1\trianglerighteq\cdots\trianglerighteq G_r=1
\]
is a normal series with abelian factors.  Thus \(G_0\) is solvable.
\end{proof}

\begin{remark}
The classical theorem in characteristic zero says that a polynomial is
solvable by radicals if and only if its Galois group is solvable.  The converse
requires Kummer theory, which is outside these notes.  The proposition above
captures the group-theoretic mechanism after a radical tower has been placed
inside a suitable Galois tower; it does not prove the converse.  Accordingly,
we do not claim a proof here of the general quintic theorem.
\end{remark}

\begin{definition}
A real number is \emph{constructible over \(\mathbb Q\)} if it belongs to a
tower of subfields of \(\mathbb R\)
\[
\mathbb Q=K_0\subseteq K_1\subseteq\cdots\subseteq K_r\subseteq\mathbb R
\]
whose successive degrees are at most two.  This is the field-theoretic form
of straightedge-and-compass
construction: line intersections use field operations and circle
intersections require a quadratic equation.
\end{definition}

The algebraic definition records a necessary consequence of the ordinary
geometric construction rules.  Indeed, the intersection of two lines has
coordinates obtained by field operations from the preceding coordinates.
After eliminating one variable, the intersection of a line and a circle, or
of two circles, is obtained by solving a quadratic equation.  Hence every
ordinary straightedge-and-compass construction produces coordinates in a
tower whose successive degrees are at most two.  This one direction is all
that is needed for the impossibility results below.

\begin{proposition}[Degree obstruction]\label{prop:construction-degree}
If \(\alpha\) is constructible over \(\mathbb Q\), then
\[
[\mathbb Q(\alpha):\mathbb Q]
\]
divides a power of \(2\).
\end{proposition}
\begin{proof}
The construction operations give a tower with successive degrees at most two.
The Tower Law makes the degree of its top field a power of two, and the degree
of the subfield \(\mathbb Q(\alpha)\) divides that degree.
\end{proof}

\begin{corollary}
Doubling the cube and trisecting a general angle are impossible with
straightedge and compass.
\end{corollary}
\begin{proof}
\(\sqrt[3]2\) has irreducible minimal polynomial \(X^3-2\).  For angle
trisection, a \(60^\circ\) angle is constructible.  If arbitrary angle
trisection were possible, then \(20^\circ\) would be constructible.  But
\(\cos20^\circ\) satisfies
\[
8X^3-6X-1=0,
\]
which is irreducible over \(\mathbb Q\) by the rational-root test.  Both
degrees are \(3\), so Proposition~\ref{prop:construction-degree} applies.
Squaring the circle uses the transcendence of \(\pi\), outside the present
scope.
\end{proof}

\section{Direct and Inverse Limits}

Infinite Galois theory is assembled from finite stages.  Direct limits join
objects that grow; inverse limits record compatible information at all finite
quotients.  The common indexing language is designed so that a finite number
of comparisons can always be made at one later stage.

\begin{definition}
A partially ordered set \(I\) is \emph{directed} if it is nonempty and every
pair \(i,j\) has an upper bound \(k\).  Examples are \(\mathbb N\), finite
subsets of a set ordered by inclusion, and finite intermediate extensions
ordered by inclusion.  Repeatedly applying the pair condition gives a common
upper bound for every finite collection.  This is exactly what lets a direct
limit turn a finite sum of representatives into one representative later on.
\end{definition}

\begin{definition}
A \emph{direct system} of \(R\)-modules over \(I\) is a family
\((M_i)_{i\in I}\) and maps
\[
\phi_i^j:M_i\longrightarrow M_j\qquad(i\leq j)
\]
with \(\phi_i^i=1_{M_i}\) and
\[
\phi_j^k\phi_i^j=\phi_i^k
\qquad(i\leq j\leq k).
\]
For \(i\leq j\leq k\), this says that the following triangle commutes:
\[
\begin{tikzcd}[column sep=large,row sep=large]
M_i \arrow[rr,"\phi_i^k"] \arrow[dr,"\phi_i^j"'] && M_k \\
& M_j \arrow[ur,"\phi_j^k"'] &
\end{tikzcd}
\]
Thus passing from \(i\) to \(k\) does not depend on whether one pauses at
\(j\).
\end{definition}

\Needspace{24\baselineskip}
\begin{definition}
A \emph{direct limit} of a direct system is an \(R\)-module
\(\varinjlim_iM_i\), together with maps
\[
\iota_i:M_i\longrightarrow\varinjlim_iM_i,
\]
such that \(\iota_j\phi_i^j=\iota_i\) for \(i\leq j\), and satisfying the
following universal property: whenever maps \(f_i:M_i\to N\) satisfy
\[
f_j\phi_i^j=f_i,
\]
there is a unique map
\[
u:\varinjlim_iM_i\longrightarrow N
\]
with \(u\iota_i=f_i\) for every \(i\).
\end{definition}

\[
\begin{tikzcd}[column sep=large,row sep=large]
M_i \arrow[rr,"\phi_i^j"] \arrow[dr,"\iota_i"'] \arrow[ddr,bend right=35,"f_i"']&&
M_j \arrow[dl,"\iota_j"] \arrow[ddl,bend left=35,"f_j"]\\
&\varinjlim_iM_i \arrow[d,dashed,"u"]\\
&N
\end{tikzcd}
\]

\begin{proposition}[Uniqueness of direct limits]
Two direct limits of the same direct system are uniquely isomorphic by an
isomorphism commuting with every canonical map.
\end{proposition}
\begin{proof}
Let \((M,\iota_i)\) and \((M',\iota_i')\) be direct limits.  The universal
property of \(M\), applied to the compatible family \((\iota_i')\), gives a
unique \(u:M\to M'\) with \(u\iota_i=\iota_i'\).  Similarly there is a unique
\(v:M'\to M\) with \(v\iota_i'=\iota_i\).  Both \(vu\) and \(1_M\) have the
same composites with every \(\iota_i\), so uniqueness gives \(vu=1_M\).
Likewise \(uv=1_{M'}\).
\end{proof}

\begin{proposition}[Construction of direct limits]
Every direct system of \(R\)-modules has a direct limit.  It is
\[
\varinjlim_{i\in I}M_i=
\left(\bigoplus_{i\in I}M_i\right)\bigg/D,
\]
where \(D\) is generated by
\[
e_i(m)-e_j(\phi_i^j(m))\qquad(i\leq j,\ m\in M_i).
\]
\end{proposition}
\begin{proof}
Let \(\iota_i(m)\) be the class of \(e_i(m)\).  The displayed relations give
\(\iota_j\phi_i^j=\iota_i\).  A compatible family \(f_i:M_i\to N\) defines
\[
\sum_i e_i(m_i)\longmapsto\sum_i f_i(m_i).
\]
The sums are finite because the source is a direct sum, and each defining
relation maps to
\[
f_i(m)-f_j(\phi_i^j(m))=0.
\]
Hence the map descends to a map \(u\) from the quotient.  It satisfies
\(u\iota_i=f_i\).  Since every quotient class is a finite sum of classes from
the summands, these identities also force uniqueness.
\end{proof}

\begin{proposition}[Maps induced on direct limits]\label{prop:direct-limit-maps}
Suppose \((M_i,\phi_i^j)\) and \((N_i,\psi_i^j)\) are direct systems and maps
\(\rho_i:M_i\to N_i\) satisfy
\[
\rho_j\phi_i^j=\psi_i^j\rho_i
\qquad(i\leq j).
\]
Then there is a unique map
\[
\rho:\varinjlim_iM_i\longrightarrow\varinjlim_iN_i
\]
with \(\rho\iota_i=\kappa_i\rho_i\), where \(\kappa_i\) are the canonical maps
for the second limit.
\end{proposition}
\begin{proof}
The maps \(\kappa_i\rho_i\) are compatible, since
\[
\kappa_j\rho_j\phi_i^j
=\kappa_j\psi_i^j\rho_i
=\kappa_i\rho_i.
\]
The universal property of \(\varinjlim_iM_i\) gives exactly the asserted map.
\end{proof}

\begin{proposition}[Element criteria]\label{prop:direct-limit-elements}
Every element of a direct limit has the form \(\iota_i(m)\) for some
\(i\in I\) and \(m\in M_i\).  In addition,
\[
\iota_i(m)=0
\quad\Longleftrightarrow\quad
\phi_i^j(m)=0\text{ for some }j\geq i,
\]
and
\[
\iota_i(m)=\iota_j(n)
\quad\Longleftrightarrow\quad
\phi_i^k(m)=\phi_j^k(n)\text{ for some }k\geq i,j.
\]
\end{proposition}
\begin{proof}
A quotient element is represented by
\[
\sum_{\ell=1}^re_{i_\ell}(m_\ell).
\]
Choose \(k\) above every \(i_\ell\).
Compatibility gives
\[
\sum_{\ell=1}^r\iota_{i_\ell}(m_\ell)
=\iota_k\left(\sum_{\ell=1}^r\phi_{i_\ell}^k(m_\ell)\right),
\]
which proves the first assertion.  The reverse implications in the displayed
criteria are compatibility.  For the forward implications, the asserted
zero or equality in the quotient is a finite sum of defining relations.
Choose an index above all of the finitely many stages occurring in this
expression.  Mapping every summand to that stage cancels the relations and
gives the desired zero or equality.
\end{proof}

\begin{example}[A directed union written as fractions]
Consider the direct system indexed by \(n\geq0\)
\[
\mathbb Z\xrightarrow{\times2}\mathbb Z\xrightarrow{\times2}\mathbb Z
\xrightarrow{\times2}\cdots.
\]
Its direct limit is naturally
\[
\mathbb Z[1/2]
=\left\{\frac{m}{2^n}:m\in\mathbb Z,\ n\geq0\right\}.
\]
Indeed, send the integer \(m\) at stage \(n\) to \(m/2^n\).  The transition
map sends \(m\) to \(2m\), and
\[
\frac{2m}{2^{n+1}}=\frac{m}{2^n},
\]
so these maps are compatible.  They induce a map from the direct limit to
\(\mathbb Z[1/2]\), and it is onto by definition.  If two representatives
\(m\) at stage \(n\) and \(a\) at stage \(r\) give the same fraction, then
after multiplying their equality by a sufficiently large power of \(2\),
their images agree at a common later stage.  Proposition~\ref{prop:direct-limit-elements}
therefore makes the induced map one-to-one.  This example turns the abstract
identification of forward representatives into the familiar rule for
equivalent fractions.
\end{example}

\begin{corollary}[Directed unions]
If \(M_i\subseteq M_j\) for \(i\leq j\) and the transition maps are
inclusions, then
\[
\varinjlim_iM_i\cong\bigcup_iM_i.
\]
\end{corollary}
\begin{proof}
The inclusions into the union are compatible.  The induced map is onto, and
the zero criterion in Proposition~\ref{prop:direct-limit-elements} makes it
one-to-one.
\end{proof}

\begin{theorem}[Exactness of directed direct limits]\label{thm:direct-limit-exact}
Suppose that
\[
A_i\xrightarrow{u_i}B_i\xrightarrow{v_i}C_i
\]
is exact for every \(i\), and all transition maps commute with these maps.
Then
\[
\varinjlim_iA_i\longrightarrow\varinjlim_iB_i
\longrightarrow\varinjlim_iC_i
\]
is exact.  Thus directed direct limits preserve injections and surjections
when these occur stagewise.
\end{theorem}
\begin{proof}
The composite is zero at every stage.  Conversely, suppose that
\(\iota_i(b)\) maps to zero.  The zero criterion gives \(j\geq i\) with
\[
v_j(\phi_i^j(b))=0.
\]
At stage \(j\), exactness supplies \(a\in A_j\) with
\(u_j(a)=\phi_i^j(b)\).  Therefore \(\iota_i(b)\) is the image of
\(\iota_j(a)\).  Applying this to sequences beginning or ending in zero gives
the final statement.
\end{proof}

\begin{definition}
An \emph{inverse system} of \(R\)-modules over the directed set \(I\) is a
family \((N_i)_{i\in I}\) and maps
\[
\psi_i^j:N_j\longrightarrow N_i\qquad(i\leq j)
\]
such that \(\psi_i^i=1_{N_i}\) and
\[
\psi_i^j\psi_j^k=\psi_i^k
\qquad(i\leq j\leq k).
\]
For \(i\leq j\leq k\), the corresponding reversed triangle commutes:
\[
\begin{tikzcd}[column sep=large,row sep=large]
& N_j \arrow[dr,"\psi_i^j"] & \\
N_k \arrow[ur,"\psi_j^k"] \arrow[rr,"\psi_i^k"'] && N_i
\end{tikzcd}
\]
Thus information at a later stage restricts consistently to every earlier
stage.  Unlike a direct limit, an inverse limit will retain only collections
of data that agree under all these restrictions.
\end{definition}

\begin{definition}
An \emph{inverse limit} of an inverse system is an \(R\)-module
\[
\varprojlim_iN_i,
\]
together with maps \(\pi_i:\varprojlim_iN_i\to N_i\) satisfying
\[
\psi_i^j\pi_j=\pi_i
\qquad(i\leq j),
\]
with the following universal property.  Whenever maps \(g_i:X\to N_i\)
satisfy \(\psi_i^jg_j=g_i\), there is a unique map
\[
u:X\longrightarrow\varprojlim_iN_i
\]
for which \(\pi_i u=g_i\) for every \(i\).
\end{definition}

\[
\begin{tikzcd}[column sep=large,row sep=large]
& X \arrow[d,dashed,"u"] \arrow[ddl,bend right=35,"g_i"']
  \arrow[ddr,bend left=35,"g_j"] & \\
& \varprojlim_iN_i \arrow[dl,"\pi_i"] \arrow[dr,"\pi_j"'] & \\
N_i && N_j \arrow[ll,"\psi_i^j"]
\end{tikzcd}
\]

\begin{proposition}[Uniqueness of inverse limits]
Two inverse limits of the same inverse system are uniquely isomorphic by an
isomorphism commuting with all projections.
\end{proposition}
\begin{proof}
Let \((N,\pi_i)\) and \((N',\pi_i')\) be inverse limits.  The universal
property of \(N'\), applied to the compatible family \((\pi_i)\), produces a
unique \(u:N\to N'\) with \(\pi_i'u=\pi_i\).  Similarly there is a unique
\(v:N'\to N\) with \(\pi_iv=\pi_i'\).  The maps \(vu\) and \(1_N\) have the
same composites with every \(\pi_i\); uniqueness gives \(vu=1_N\).  The
same argument gives \(uv=1_{N'}\).
\end{proof}

\begin{proposition}[Construction of inverse limits]
Every inverse system of \(R\)-modules has an inverse limit.  It is
\[
\varprojlim_{i\in I}N_i=
\left\{(n_i)\in\prod_{i\in I}N_i:
\psi_i^j(n_j)=n_i\text{ whenever }i\leq j\right\}.
\]
\end{proposition}
\begin{proof}
The displayed set is a submodule of the product: the compatibility equations
are preserved under addition and scalar multiplication.  Its coordinate
projections are compatible.  A compatible family \(g_i:X\to N_i\) gives the
map
\[
x\longmapsto(g_i(x))_{i\in I}
\]
to this submodule, because the assumed equations say precisely that this
tuple is compatible.  Its \(i\)-th coordinate is \(g_i(x)\), and a map to a
submodule of a product is determined by its coordinates.  This proves the
universal property.  Thus a direct limit is a quotient of a direct sum,
whereas an inverse limit is a compatible submodule of a direct product.
\end{proof}

\begin{proposition}[Maps induced on inverse limits]
Suppose \((M_i,\phi_i^j)\) and \((N_i,\psi_i^j)\) are inverse systems and
maps \(\rho_i:M_i\to N_i\) satisfy
\[
\psi_i^j\rho_j=\rho_i\phi_i^j
\qquad(i\leq j).
\]
Then there is a unique map
\[
\rho:\varprojlim_iM_i\longrightarrow\varprojlim_iN_i
\]
whose \(i\)-th coordinate is \(\rho_i\) applied to the \(i\)-th coordinate.
\end{proposition}
\begin{proof}
For \((m_i)_i\in\varprojlim_iM_i\), the compatibility calculation
\[
\psi_i^j\bigl(\rho_j(m_j)\bigr)
=\rho_i\bigl(\phi_i^j(m_j)\bigr)
=\rho_i(m_i)
\]
shows that \((\rho_i(m_i))_i\) is a compatible tuple.  The asserted formula
therefore defines \(\rho\); uniqueness follows from the coordinate
projections.
\end{proof}

\begin{example}[The \(2\)-adic and profinite integers]
The inverse limit of the reductions
\[
\mathbb Z/2^{n+1}\mathbb Z\longrightarrow\mathbb Z/2^n\mathbb Z
\]
is \(\mathbb Z_2\).  An element is a compatible sequence
\[
(a_n\bmod2^n)_{n\geq1};
\]
the ordinary integer \(5\), for instance, gives
\((1\bmod2,1\bmod4,5\bmod8,5\bmod16,\ldots)\).  Ordered by divisibility,
\[
\widehat{\mathbb Z}=\varprojlim_{n\geq1}\mathbb Z/n\mathbb Z.
\]
The former retains only powers of \(2\); the latter retains every modulus.
\(\mathbb Z_2\) also contains compatible sequences not coming from a
nonnegative ordinary integer in any fixed finite stage.  For example,
\[
(-1\bmod2,\,-1\bmod4,\,-1\bmod8,\,-1\bmod16,\ldots)
\]
is compatible and represents \(-1\), while the proposed first two entries
\[
(0\bmod2,\ 1\bmod4)
\]
cannot occur in a compatible sequence, because \(1\bmod4\) reduces to
\(1\bmod2\), not to \(0\bmod2\).  Compatibility is therefore a real
restriction, not merely a way of listing residues.
\end{example}

\begin{proposition}[Left exactness of inverse limits]
Inverse limits preserve kernels and hence are left exact, but need not
preserve surjectivity at the right end.
\end{proposition}
\begin{proof}
In the product description, a compatible family lies in the kernel of an
induced map exactly when every coordinate lies in the corresponding kernel.
The same compatibility equations remain, so the kernel is the inverse limit
of the kernels.  Compatible lifts of a compatible family need not exist,
which is why right exactness is not claimed.
\end{proof}

\begin{example}[Stagewise surjections need not lift compatibly]
For every \(n\geq1\), let \(X_n=\mathbb Z\) with identity transition maps,
let
\[
Y_n=\mathbb Z/2^n\mathbb Z
\]
with the reduction maps, and let \(q_n:X_n\to Y_n\) be the quotient map.
Each \(q_n\) is surjective, and the maps commute with the transitions.  Yet
\[
\varprojlim_nX_n\cong\mathbb Z,
\qquad
\varprojlim_nY_n\cong\mathbb Z_2,
\]
and the induced map \(\mathbb Z\to\mathbb Z_2\) is not surjective.  For
example, the compatible sequence
\[
(1\bmod2,\ 1\bmod4,\ 5\bmod8,\ 5\bmod16,\ 21\bmod32,\ldots)
\]
is not the image of an ordinary integer.  Its \(2\)-adic binary expansion has
infinitely many zeros and ones.  A nonnegative ordinary integer has only
finitely many nonzero binary digits, whereas a negative ordinary integer has a
\(2\)-adic binary expansion that is eventually all ones.  Thus compatible
stagewise lifts need not be selectable simultaneously.
\end{example}

\begin{remark}[Looking Ahead: the Mittag--Leffler condition]
The preceding proposition identifies the difficulty with inverse limits: one
can take compatible kernels coordinatewise, but compatible choices of
preimages need not exist.  For an inverse system \((M_i,\psi_i^j)\), the
\emph{Mittag--Leffler condition} says that for every fixed \(i\), there is a
\(j\geq i\) such that
\[
\operatorname{im}(\psi_i^k)=\operatorname{im}(\psi_i^j)
\qquad\text{for every }k\geq j.
\]
In other words, the images inside the fixed earlier module \(M_i\), as seen
from later and later stages, eventually stabilize.  Surjective transition
maps satisfy the condition because all these images are \(M_i\).  It also
holds when each \(M_i\) is finite as a set, since a descending family of
submodules of a fixed finite module stabilizes; for finite-dimensional vector
spaces, one may instead use the eventually constant nonincreasing dimensions
of the images.

Under standard hypotheses, a Mittag--Leffler condition on the kernel system
in an inverse system of short exact sequences restores surjectivity at the
right-hand end after taking inverse limits.  It is not needed for the present
Galois application: restriction maps between finite Galois stages are already
surjective.  The derived obstruction to exactness is usually denoted
\(\varprojlim{}^{1}\); its systematic treatment belongs to later homological
algebra.
\end{remark}

\section{Topological Preliminaries}

In finite Galois theory the group alone remembers all the relevant
information.  For an infinite extension, however, one must also remember when
automorphisms agree on larger and larger finite Galois subextensions.  Topology
is the concise language for this finite-stage approximation: in the Krull
topology, two automorphisms are close when they agree on a sufficiently large
finite Galois field.  We develop only the small amount of topology needed to
make this statement precise.

\subsection{Basic Point-Set Topology}

\begin{definition}
A \emph{topology} on a set \(X\) is a collection \(\tau\) of subsets of \(X\)
such that
\[
\varnothing,\ X\in\tau;
\]
an arbitrary union of members of \(\tau\) is in \(\tau\); and a finite
intersection of members of \(\tau\) is in \(\tau\).  The members of \(\tau\)
are the \emph{open sets}.  A set is \emph{closed} if its complement is open,
and a \emph{neighborhood} of \(x\in X\) is a set containing an open set that
contains \(x\).
\end{definition}

The discrete topology, in which every subset is open, is the important
example here: every finite Galois group will be given this topology.  At the
opposite extreme, the indiscrete topology has only \(\varnothing\) and \(X\)
open.  The usual topology on \(\mathbb R\) is familiar intuition, but no
analysis on \(\mathbb R\) is needed below.

\begin{definition}
A collection \(\mathcal B\) of open sets is a \emph{basis} if every open set
is a union of members of \(\mathcal B\).  A collection of neighborhoods of
\(x\) is a \emph{neighborhood basis at \(x\)} if every neighborhood of \(x\)
contains one of its members.  A map \(f:X\to Y\) is \emph{continuous} if
\(f^{-1}(U)\) is open whenever \(U\) is open.  A bijective continuous map
whose inverse is continuous is a \emph{homeomorphism}.
\end{definition}

For \(A\subseteq X\), its \emph{closure} \(\overline A\) is the intersection
of all closed subsets of \(X\) containing \(A\).  We say that \(A\) is
\emph{dense} when \(\overline A=X\).  For example, a nonempty subset of a
finite discrete space is closed, so it is dense exactly when it is the whole
space.

If \(Y\subseteq X\), the \emph{subspace topology} on \(Y\) declares
\[
U\subseteq Y\text{ open}\quad\Longleftrightarrow\quad
U=Y\cap V\text{ for some open }V\subseteq X.
\]
For a family of spaces, the product topology on
\[
\prod_{i\in I}X_i
\]
has as a basis the sets which restrict only finitely many coordinates to open
sets and leave every remaining coordinate unrestricted.  Thus a basic
neighborhood records finite information.  When every \(X_i\) is finite
discrete, it may simply prescribe the values of finitely many coordinates.
For example, in
\[
\prod_{n\geq1}\mathbb Z/2^n\mathbb Z,
\]
the condition
\[
a_1=1\bmod2,
\qquad
a_2=3\bmod4
\]
with all higher coordinates unrestricted is a basic open set of the full
product.  Compatibility is imposed only after one passes to the inverse-limit
subspace.
Inverse limits will be compatible subspaces of precisely such products.

\begin{definition}
A space is \emph{Hausdorff} if distinct points have disjoint neighborhoods.
It is \emph{compact} if every open cover has a finite subcover.  A space is
\emph{connected} if it cannot be written as the union of two disjoint nonempty
open sets, and \emph{totally disconnected} if every connected subset has at
most one point.
\end{definition}

Finite discrete spaces are Hausdorff, compact, and totally disconnected:
singletons separate points and are open, while an open cover of a finite set
has an evident finite subcover.

\begin{theorem}[Tychonoff's theorem {\rm(external)}]
An arbitrary product of compact spaces is compact in the product topology.
\end{theorem}

\begin{remark}[Foundational scope]
In standard set theory with the Axiom of Choice, Tychonoff's theorem gives
compactness for arbitrary products of compact spaces.  In fact, over ZF, this
full arbitrary-product form is equivalent to the Axiom of Choice.  The present
notes use it only as an external input for products of finite discrete spaces;
this is consistent with the earlier use of Zorn's Lemma in the construction of
algebraic closures.
\end{remark}

\subsection{Topological Groups}

\begin{definition}
A \emph{topological group} is a group \(G\) with a topology for which
\[
G\times G\longrightarrow G,\qquad (x,y)\longmapsto xy,
\]
and
\[
G\longrightarrow G,\qquad x\longmapsto x^{-1},
\]
are continuous.  A homomorphism of topological groups is \emph{continuous}
when it is continuous as a map of spaces; a \emph{topological-group
isomorphism} is a group isomorphism that is also a homeomorphism.
\end{definition}

\begin{proposition}
For \(g\in G\), left and right translation,
\[
L_g(x)=gx,\qquad R_g(x)=xg,
\]
are homeomorphisms of a topological group.  Consequently, the neighborhoods
of the identity determine all neighborhoods: if \(U\) is a neighborhood of
\(1\), then \(gU\) is a neighborhood of \(g\).
\end{proposition}
\begin{proof}
The map \(x\mapsto(g,x)\) is continuous, and left translation is its
composite with multiplication.  Similarly, \(x\mapsto(x,g)\) followed by
multiplication is right translation.  Their inverse translations are
\(L_{g^{-1}}\) and \(R_{g^{-1}}\), which are continuous by the same argument.
The final assertion follows by applying \(L_g\) to a neighborhood of \(1\).
\end{proof}

\begin{proposition}
If \(\phi:G\to H\) is a continuous homomorphism and \(H\) is discrete, then
\(\ker\phi\) is open.  If \(H\) is finite, this kernel has finite index.
Further, every subgroup with its subspace topology and every product of
topological groups with coordinatewise operations are topological groups.
\end{proposition}
\begin{proof}
Since \(\{1\}\) is open in a discrete group,
\[
\ker\phi=\phi^{-1}(\{1\})
\]
is open.  The First Isomorphism Theorem identifies \(G/\ker\phi\) with the
finite group \(\operatorname{im}\phi\).  The remaining assertions follow by
restricting the continuous group operations, respectively by checking
continuity coordinatewise on basic product opens.
\end{proof}

\subsection{Profinite Groups}

\begin{definition}
A \emph{profinite group} is an inverse limit of finite discrete groups, with
the subspace topology inherited from
\[
\prod_{i\in I}G_i.
\]
The topology says that two elements are close when they have the same image at
many finite stages.  The next proposition makes this intuition precise in the
particularly useful form of a neighborhood basis.
\end{definition}

One may think of a profinite group as an infinite symmetry group observed
through compatible finite stages.  No single stage sees the entire element,
but a compatible answer at every stage does.
The prototype is \(\widehat{\mathbb Z}\): an element gives a residue class
modulo every positive integer, with all reductions agreeing.  This viewpoint
is exactly suited to an infinite algebraic extension, whose individual
elements and automorphisms can be tested on finite subextensions.

\begin{proposition}
An inverse limit of finite discrete groups is a closed topological subgroup of
their product.  Consequently every profinite group is compact, Hausdorff, and
totally disconnected.
\end{proposition}
\begin{proof}
The product is a topological group and is compact by Tychonoff's theorem.  It
is Hausdorff: distinct points differ in one coordinate, and disjoint open
neighborhoods in that finite discrete factor pull back to disjoint product
neighborhoods.  It is totally disconnected for the same clopen-cylinder
reason: if \(x\ne y\), choose a coordinate where they differ; fixing that
coordinate to \(x_i\) gives a clopen set containing \(x\) but not \(y\), so
no connected subset can contain both points.  For \(i\leq j\), define
\[
\rho_i^j:\prod_kG_k\longrightarrow G_i\times G_i,
\qquad
(g_k)_k\longmapsto\bigl(\psi_i^j(g_j),g_i\bigr).
\]
This map is continuous.  The compatibility condition
\(\psi_i^j(g_j)=g_i\) is the inverse image under \(\rho_i^j\) of the
diagonal, which is closed because the finite discrete target is Hausdorff.
Intersecting these closed conditions over all \(i\leq j\) gives the inverse
limit.  A closed subspace of a compact space is compact, since an open cover
of it extends by its open complement to an open cover of the ambient space.
Hausdorffness and total disconnectedness pass to subspaces: separating open
sets and clopen cylinders may simply be intersected with the subspace.
\end{proof}

\begin{proposition}[The projection-kernel basis]\label{prop:profinite-open-subgroups}
Let
\[
G=\varprojlim_{i\in I}G_i
\]
be a profinite group, indexed by a directed set.  The subgroups
\[
\ker(\pi_i)\qquad(i\in I)
\]
form a neighborhood basis at the identity.  They are open normal subgroups,
and every open subgroup has finite index and is closed.
\end{proposition}
\begin{proof}
In the product topology, a basic neighborhood of the identity restricts only
finitely many coordinates, say those indexed by \(i_1,\ldots,i_r\).  Choose
\(k\) above all these indices.  Compatibility implies
\[
\ker(\pi_k)\subseteq
\bigcap_{\ell=1}^r\ker(\pi_{i_\ell}),
\]
so one projection kernel lies inside the given neighborhood.  Conversely,
each \(\ker(\pi_i)\) is open because it is the inverse image of the open
singleton \(\{1\}\subseteq G_i\).  It is normal, and its quotient is a
subgroup of the finite group \(G_i\).  An open subgroup contains one such
kernel, hence has finite index.  Its complement is a finite union of its open
cosets, so it is closed.
\end{proof}

\begin{example}[The topology of \(\mathbb Z_p\)]
For a prime \(p\), define
\[
\mathbb Z_p=\varprojlim_{n\geq1}\mathbb Z/p^n\mathbb Z
\]
using the usual reduction maps.  The projection
\[
\mathbb Z_p\longrightarrow\mathbb Z/p^n\mathbb Z
\]
has kernel customarily denoted \(p^n\mathbb Z_p\).  Thus
\[
p^n\mathbb Z_p\qquad(n\geq1)
\]
is a neighborhood basis at \(0\).  In particular, congruence modulo a high
power of \(p\) is the topological notion of being close in \(\mathbb Z_p\).
This is the additive prototype for the Krull topology below.
\end{example}

\section{Infinite Galois Theory}

Let \(E/F\) be algebraic Galois.  Each element of \(E\) lies in a finite
Galois subextension: take the splitting field in \(E\) of its separable
minimal polynomial.  If \(K/F\) and \(L/F\) are finite Galois subextensions,
their \emph{compositum} \(KL\) is the smallest subfield of \(E\) containing
both.  It is finite over \(F\), separable because it lies in the separable
extension \(E/F\), and normal because \(K\) and \(L\) are splitting fields
over \(F\) of finitely many separable polynomials, while \(KL\) is the
splitting field of the product of those polynomials.  Thus \(KL/F\) is again
finite Galois and contains both fields.  These subextensions are therefore
directed by compositum, and
\[
E=\bigcup_{K/F\ \mathrm{finite\ Galois}}K.
\]
For the running finite-field example, the finite stages inside
\(\mathbb F_q^{\mathrm{sep}}\) are the fields \(\mathbb F_{q^n}\), ordered
by divisibility of \(n\).  Restriction from a larger finite field to a smaller
one is the concrete model for every inverse-limit transition map used below.

\begin{lemma}[Extension to an algebraic closure]\label{lem:closure-auto-extension}
Let \(L/F\) be algebraic inside a fixed algebraic closure \(\overline F\).
Every \(F\)-embedding
\[
\sigma:L\longrightarrow\overline F
\]
extends to an element of \(\operatorname{Aut}_F(\overline F)\).
\end{lemma}
\begin{proof}
The embedding-extension theorem (Theorem~\ref{thm:embedding-extension}) extends \(\sigma\) to an \(F\)-embedding
\[
\widetilde\sigma:\overline F\longrightarrow\overline F.
\]
Its image is an algebraically closed subfield.  Since \(\overline F/F\) is
algebraic, \(\overline F/\widetilde\sigma(\overline F)\) is algebraic.  An
algebraically closed field has no proper algebraic extension: if \(B/A\) is
algebraic and \(A\) is algebraically closed, then the minimal polynomial over
\(A\) of any \(\beta\in B\) has a root in \(A\), so irreducibility forces it
to be linear.  Hence \(\beta\in A\).  Thus
\(\widetilde\sigma\) is onto, as required.
\end{proof}

\begin{lemma}[Extension through a normal algebraic extension]
If \(E/F\) is algebraic and normal and \(L/F\) is an algebraic subextension
of \(E/F\), every \(F\)-embedding \(L\to\overline F\) that maps \(L\) into
\(E\) extends to an \(F\)-automorphism of \(E\).  In particular, every
\(F\)-automorphism of a finite subextension \(L/F\) extends to an
\(F\)-automorphism of \(E\).
\end{lemma}
\begin{proof}
Extend the given embedding to
\(\widetilde\sigma\in\operatorname{Aut}_F(\overline F)\) by
Lemma~\ref{lem:closure-auto-extension}.  If \(x\in E\), then
\(\widetilde\sigma(x)\) is another root of the minimal polynomial of \(x\)
over \(F\).  Normality says that all such roots lie in \(E\), so
\(\widetilde\sigma(E)\subseteq E\).  Applying the same argument to
\(\widetilde\sigma^{-1}\) gives the reverse inclusion.  Hence
\(\widetilde\sigma|_E\) is the desired automorphism.
\end{proof}

\begin{theorem}[Infinite Galois groups as inverse limits]\label{thm:infinite-galois-limit}
For a Galois extension \(E/F\), restriction induces an isomorphism
\[
\operatorname{Gal}(E/F)\cong
\varprojlim_{K/F\ \mathrm{finite\ Galois}}\operatorname{Gal}(K/F).
\]
\end{theorem}
\begin{proof}
For \(K\subseteq L\), restriction
\(\operatorname{Gal}(L/F)\to\operatorname{Gal}(K/F)\) is well-defined and
surjective by extension through the finite normal extension \(L/F\), and the
composition law is immediate.

An automorphism of \(E\) gives compatible restrictions.  It is determined by
them because the finite Galois fields cover \(E\).  Conversely, for a
compatible family \((\sigma_K)_K\), define \(\sigma(x)=\sigma_K(x)\) at any
finite Galois stage containing \(x\).  A common larger stage proves this is
well-defined and respects addition and multiplication.  It fixes \(F\), and
the compatible inverse family defines its inverse.  Thus it is an
automorphism of \(E\).
\end{proof}

\begin{definition}
The \emph{Krull topology} on \(\operatorname{Gal}(E/F)\) has basic
neighborhoods of the identity
\[
\operatorname{Gal}(E/K),
\]
where \(K/F\) is finite Galois.  Their cosets are the other basic open sets.
\end{definition}

This topology is forced by the finite-stage viewpoint: two automorphisms
\(\sigma,\tau\) are close when they agree on a large finite Galois
subextension \(K/F\).  Equivalently,
\[
\tau^{-1}\sigma\in\operatorname{Gal}(E/K).
\]
Thus the displayed stabilizers are exactly the natural identity
neighborhoods, and their translates are all other basic neighborhoods.
They do form a basis: for finite Galois \(K_1/F\) and \(K_2/F\), the
directedness argument above gives a finite Galois compositum \(K_1K_2/F\),
and
\[
\operatorname{Gal}(E/K_1)\cap\operatorname{Gal}(E/K_2)
=\operatorname{Gal}(E/K_1K_2).
\]

Under Theorem~\ref{thm:infinite-galois-limit}, these subgroups are exactly the
kernels of finite-stage projections.  Thus the Krull topology is the
inverse-limit topology, the displayed isomorphism is an isomorphism of
topological groups, and \(\operatorname{Gal}(E/F)\) is profinite.
Concretely, in the finite-field case a basic neighborhood of the identity
consists of the automorphisms that fix a prescribed
\(\mathbb F_{q^n}\).  The topology therefore says that two automorphisms are
close when their Frobenius powers agree on a sufficiently large finite field.

The infinite correspondence follows the same finite-stage logic as the
finite theorem, with one new feature: closedness is what guarantees that a
subgroup is recovered from all of its finite-stage images.  The proof
first recovers an intermediate field from its stabilizer, then uses the
finite correspondence at every finite Galois stage, and finally translates
finiteness and normality into openness and normality of the subgroup.
In the finite case every subgroup is automatically closed.  In the infinite
case a nonclosed subgroup can omit automorphisms that its elements approximate
at every finite stage, so it cannot be recovered from its finite images.

\begin{theorem}[Infinite Galois correspondence]\label{thm:infinite-galois}
For a Galois extension \(E/F\), the assignments
\[
K\longmapsto\operatorname{Gal}(E/K),
\qquad
H\longmapsto E^H
\]
are inverse inclusion-reversing bijections between intermediate fields and
closed subgroups of \(\operatorname{Gal}(E/F)\).  Further,
\[
K/F\text{ finite}\quad\Longleftrightarrow\quad
\operatorname{Gal}(E/K)\text{ open},
\]
and
\[
K/F\text{ Galois}\quad\Longleftrightarrow\quad
\operatorname{Gal}(E/K)\text{ normal}.
\]
Consequently, finite Galois intermediate extensions correspond to open normal
subgroups.
\end{theorem}
\begin{proof}
For \(\alpha\in E\), choose a finite Galois subextension \(L/F\) containing
\(\alpha\).  Its point stabilizer contains the open subgroup
\(\operatorname{Gal}(E/L)\), hence is itself open.  It has finite index and
is closed by Proposition~\ref{prop:profinite-open-subgroups}.  Therefore, for
an intermediate \(K\),
\[
\operatorname{Gal}(E/K)=\bigcap_{\alpha\in K}\operatorname{Stab}(\alpha)
\]
is closed.  We have
\[
E^{\operatorname{Gal}(E/K)}=K. \tag{12.2}
\]
Indeed, if \(\alpha\notin K\), its minimal polynomial over \(K\) has degree
larger than one.  The extension \(E/K\) is normal and separable, so that
polynomial has in \(E\) a second, distinct root \(\beta\).  The map
\(K(\alpha)\to K(\beta)\) fixing \(K\) and taking \(\alpha\) to \(\beta\)
extends, by the preceding lemma applied over \(K\), to a \(K\)-automorphism of
\(E\).  It moves \(\alpha\), proving (12.2).

Let \(H\) be closed, and let \(H_L\) be its image in
\(\operatorname{Gal}(L/F)\) for a finite Galois \(L/F\).  The field
\(L^{H_L}\) lies in \(E^H\).  If \(\sigma\) fixes \(E^H\), then
\(\sigma|_L\) fixes \(L^{H_L}\), so the finite correspondence gives
\(\sigma|_L\in H_L\).  Thus the image of \(\sigma\) lies in the image of
\(H\) at every finite stage; equivalently, every basic open neighborhood of
\(\sigma\) meets \(H\).  Hence \(\sigma\) lies in the closure of \(H\), and so
in \(H\).  Therefore \(\operatorname{Gal}(E/E^H)=H\), completing the
correspondence.

If \(K/F\) is finite, then it is separable; by the
\hyperref[thm:finite-galois-closure]{finite Galois closure theorem} and
normality of \(E/F\), choose a finite Galois extension \(L/F\) in \(E\)
containing \(K\).
Then \(\operatorname{Gal}(E/L)\) is open and is contained in
\(\operatorname{Gal}(E/K)\).  Conversely, openness supplies a finite Galois
\(L/F\) with \(\operatorname{Gal}(E/L)\) contained in
\(\operatorname{Gal}(E/K)\).  Taking fixed fields in (12.2) gives
\(K\subseteq L\).
For the normality assertion, write \(G=\operatorname{Gal}(E/F)\) and
\(H=\operatorname{Gal}(E/K)\).  First observe that
\[
\sigma H\sigma^{-1}=\operatorname{Gal}(E/\sigma(K))
\qquad(\sigma\in G).
\]
Indeed, an element on the left fixes \(\sigma(K)\), and the converse follows
by conjugating back.  If \(K/F\) is Galois, each \(\sigma\in G\) preserves
\(K\), so this formula gives \(\sigma H\sigma^{-1}=H\).  Conversely, if
\(H\) is normal, the formula and (12.2) give
\[
\sigma(K)=E^{\operatorname{Gal}(E/\sigma(K))}=E^H=K
\qquad(\sigma\in G).
\]
To prove that \(K/F\) is normal, let \(\tau:K\to\overline F\) be an
\(F\)-embedding.  By Lemma~\ref{lem:closure-auto-extension}, it extends to
an automorphism \(\widetilde\tau\) of \(\overline F\) fixing \(F\).  Since
\(E/F\) is normal, \(\widetilde\tau(E)=E\), so its restriction belongs to \(G\).
The displayed stability of \(K\) therefore yields \(\tau(K)=K\).  Thus
\(K/F\) is normal; it is separable because it is an intermediate extension
of the separable extension \(E/F\).  Hence it is Galois.
\end{proof}

\begin{definition}
Inside a fixed algebraic closure, \(F^{\mathrm{sep}}\) is the
\emph{separable closure} of \(F\), and
\[
G_F=\operatorname{Gal}(F^{\mathrm{sep}}/F)
\]
is its \emph{absolute Galois group}.
\end{definition}

\begin{theorem}[The absolute Galois group of a finite field]
Let \(\mathbb F_q^{\mathrm{sep}}\) be a separable closure of \(\mathbb F_q\).
Then
\[
G_{\mathbb F_q}
\cong
\varprojlim_{n\geq1}\mathbb Z/n\mathbb Z
=\widehat{\mathbb Z}
\]
as topological groups.  Under this identification, Frobenius
\[
\operatorname{Fr}_q:x\longmapsto x^q
\]
corresponds to \((1\bmod n)_{n\geq1}\) and is a dense topological generator.
\end{theorem}
\begin{proof}
Every element of \(\mathbb F_q^{\mathrm{sep}}\) has finite degree over
\(\mathbb F_q\), and hence lies in some \(\mathbb F_{q^n}\).  The finite
extensions \(\mathbb F_{q^n}\) are directed by divisibility, so
Theorem~\ref{thm:infinite-galois-limit} gives
\[
G_{\mathbb F_q}
\cong
\varprojlim_{n\geq1}\operatorname{Gal}(\mathbb F_{q^n}/\mathbb F_q).
\]
By Proposition~\ref{prop:frobenius-order}, the \(n\)-th group is cyclic of
order \(n\), generated by \(\operatorname{Fr}_q\).  Identify it with
\(\mathbb Z/n\mathbb Z\) by sending the class of \(a\) to
\(\operatorname{Fr}_q^a\).  If \(n\mid m\), restriction from
\(\mathbb F_{q^m}\) to \(\mathbb F_{q^n}\) sends
\(\operatorname{Fr}_q^a\) to \(\operatorname{Fr}_q^a\), which is exactly
the reduction map
\[
\mathbb Z/m\mathbb Z\longrightarrow\mathbb Z/n\mathbb Z.
\]
The resulting inverse limit is \(\widehat{\mathbb Z}\).  Frobenius maps to
the compatible family \((1\bmod n)_n\).  Its cyclic subgroup maps onto every
finite quotient \(\mathbb Z/n\mathbb Z\), so it meets every basic open coset;
the projection-kernel basis in Proposition~\ref{prop:profinite-open-subgroups}
therefore shows that it is dense.
\end{proof}

Here dense generator is deliberately topological language.  Frobenius does
not generate \(\widehat{\mathbb Z}\) as an abstract cyclic group; rather, the
ordinary cyclic subgroup generated by Frobenius is dense in
\(\widehat{\mathbb Z}\).

Under this identification,
\[
\operatorname{Gal}(\mathbb F_q^{\mathrm{sep}}/\mathbb F_{q^n})
\]
corresponds to
\[
\ker\bigl(\widehat{\mathbb Z}\longrightarrow\mathbb Z/n\mathbb Z\bigr)
=n\widehat{\mathbb Z},
\]
the unique open subgroup of index \(n\); its uniqueness also follows from the
infinite Galois correspondence together with the uniqueness of the degree-
\(n\) extension \(\mathbb F_{q^n}/\mathbb F_q\).  Thus the finite fields,
Frobenius, inverse limits, and the open-subgroup part of the infinite Galois
correspondence meet in one explicit example.

\section{Further Topics in Algebra}

This final section is orientation rather than prerequisite material.  Its aim
is to show how the four viewpoints developed in these notes become the common
language of later algebra.  The universal properties of free objects,
quotients, tensor products, and limits encountered throughout these notes are
organized systematically by category theory, whose formal development belongs
to later study.

\subsection{Groups, Actions, and Linear Algebra}

At its most elementary, a representation of a group \(G\) is a homomorphism
\[
G\longrightarrow\operatorname{GL}(V).
\]
Thus representation theory combines group theory with linear algebra.
Actions from Part~I are actions on sets; representations are the especially
structured case in which the set is a vector space.  Invariant subspaces,
irreducible representations, characters, and induction and restriction
organize this richer form of symmetry.

Normal subgroups, quotients, semidirect products, and exact sequences lead to
the study of group extensions.  Group cohomology later organizes extension
and obstruction questions.  Groups also arise outside algebra: fundamental
groups, covering spaces, and deck transformations make group actions central
in algebraic topology.  Matrix groups, Lie groups, and Lie algebras are a
further meeting point of group theory and linear algebra.

\subsection{Rings, Local Structure, and Geometry}

Part~II is the entry point to commutative algebra.  Ideals, localization,
polynomial rings, prime and maximal ideals, and modules lead onward to
Noetherian rings, primary decomposition, integral dependence, Krull
dimension, completions, and local algebra.  Localization deserves special
emphasis: it is algebraic zooming-in near a prime ideal.

Affine algebraic geometry studies spaces encoded by commutative rings.
Informally, \(\operatorname{Spec}(R)\) has prime ideals as points;
localization describes what is visible near a point, residue fields record
its scalar data, and modules supply algebraic analogues of geometric linear
data.  Scheme theory systematically combines rings, localization, modules,
tensor products, and field extensions.

\subsection{Modules, Exactness, and Multilinear Algebra}

Exact sequences, projective and injective modules, flatness, Hom, tensor
products, and direct and inverse limits form the starting vocabulary of
homological algebra.  Projective and injective resolutions, \(\operatorname{Tor}\),
\(\operatorname{Ext}\), derived functors, and spectral sequences are later
tools built from this vocabulary.  The Mittag--Leffler remark and the symbol
\(\varprojlim{}^1\) are early signals of that theory.

Tensor products, dual spaces, exterior powers, and alternating forms also
lead to differential geometry: they organize differential forms, orientation,
volume forms, and tangent and cotangent spaces.  No manifold theory is needed
to recognize this continuation.

\subsection{Fields, Arithmetic, and Galois Symmetry}

Finite extensions of \(\mathbb Q\) lead to number fields.  Their rings of
algebraic integers, ideals in Dedekind domains, factorization of prime ideals,
ramification, and decomposition and inertia groups combine the themes of
rings, localization, field extensions, and Galois groups.  The present notes
also deliberately stop before the full solvability-by-radicals theorem;
Kummer theory is a natural next step when suitable roots of unity are present.

Absolute Galois groups act on algebraic objects, and Galois cohomology later
organizes these actions in field arithmetic, descent, and number theory.
Continuous homomorphisms
\[
G_F\longrightarrow\operatorname{GL}_n(R)
\]
for suitable topological coefficient rings are Galois representations, a
major meeting point of Galois theory, representation theory, and topology.
Algebraic geometry has a parallel profinite fundamental group: finite
\'etale covers play the role of finite covering spaces.  Differential Galois
theory and covering-space Galois theory offer further variants of the same
symmetry principle.

\subsection{A Synthesis}

\[
\begin{array}{rcl}
\text{Groups} &\longrightarrow& \text{symmetry and actions},\\
\text{Rings} &\longrightarrow& \text{arithmetic and local structure},\\
\text{Modules} &\longrightarrow& \text{linearized algebra and exactness},\\
\text{Fields} &\longrightarrow& \text{algebraic equations and Galois symmetry}.
\end{array}
\]
Later subjects arise by combining these viewpoints: representation theory is
groups plus linear algebra; homological algebra is modules plus exact
sequences and Hom/tensor; algebraic geometry is commutative rings plus
localization, modules, and field extensions; algebraic number theory is rings
plus fields and Galois theory; and arithmetic geometry joins algebraic
geometry, number theory, and Galois representations.  These are conceptual
slogans, not formal identities.  The recurring methods of the notes---quotient,
localization, universal property, action, kernel and image, exact sequence,
tensor construction, passage to finite stages, and symmetry---continue to
reappear throughout advanced mathematics.

\section*{Exercises}
\subsection*{Basic}
\begin{exercise}[Basic] Prove directly that the fixed set of field
automorphisms is a field.
\end{exercise}
\begin{exercise}[Basic] Verify Artin independence for the two embeddings of
\(\mathbb Q(\sqrt2)\) in \(\mathbb C\).
\end{exercise}
\begin{exercise}[Basic] Determine the correspondence for
\(\mathbb Q(\sqrt2,\sqrt3)/\mathbb Q\).
\end{exercise}
\begin{exercise}[Basic] Compute the direct limit of
\(\mathbb Z\xrightarrow{\times3}\mathbb Z\xrightarrow{\times3}\cdots\).
\end{exercise}
\begin{exercise}[Basic] Describe the compatible sequence in \(\mathbb Z_2\)
associated with each of \(-1,0,1,2\).
\end{exercise}
\begin{exercise}[Basic] Prove that every discrete space is Hausdorff, and
describe a basic open set in a product of finite discrete spaces.
\end{exercise}

\subsection*{Core}
\begin{exercise}[Core] Reprove Lemma~\ref{lem:artin-independence} using any
chosen pair of distinct embeddings.
\end{exercise}
\begin{exercise}[Core] List all subgroup-field pairs in the splitting field
of \(X^3-2\), including their degrees.
\end{exercise}
\begin{exercise}[Core] Work out the correspondence for
\(\mathbb Q(\zeta_5)/\mathbb Q\), including normality.
\end{exercise}
\begin{exercise}[Core] Prove the equality criterion in
Proposition~\ref{prop:direct-limit-elements} from the quotient construction.
\end{exercise}
\begin{exercise}[Core] Use the element criteria to prove exactness for a
specified directed system of short exact sequences.
\end{exercise}
\begin{exercise}[Core] Give an inverse system for which surjectivity at each
stage does not give surjectivity on inverse limits.
\end{exercise}
\begin{exercise}[Core] Show that projection kernels form a neighborhood basis
in a profinite group.
\end{exercise}
\begin{exercise}[Core] Prove that the kernel of a continuous homomorphism to a
discrete group is open, and use translations to describe its cosets as open
sets.
\end{exercise}
\begin{exercise}[Core] Verify directly the gluing step in
Theorem~\ref{thm:infinite-galois-limit}.
\end{exercise}
\begin{exercise}[Core] Give a direct alternative proof that
\(\operatorname{Gal}(E/K)\) is closed in the Krull topology.
\end{exercise}
\begin{exercise}[Core] Show that an open subgroup of a profinite group has
finite index.
\end{exercise}

\subsection*{Further}
\begin{exercise}[Further] Prove the root-of-unity radical proposition in full
for a single radical step.
\end{exercise}
\begin{exercise}[Further] Use the degree obstruction to show that a regular
heptagon is not straightedge-and-compass constructible.
\end{exercise}
\begin{exercise}[Further] Show that \(\widehat{\mathbb Z}\) maps onto
\(\mathbb Z/n\mathbb Z\) for every \(n>0\).
\end{exercise}
\begin{exercise}[Further] Prove surjectivity of restriction from an infinite
Galois group to a finite Galois stage.
\end{exercise}
\begin{exercise}[Further] Let \(E/F\) be Galois and let \(K/F\) be finite.
Choose a finite Galois extension \(L/F\) in \(E\) containing \(K\).  Use the
finite correspondence in \(L/F\) to decide when
\(\operatorname{Gal}(E/K)\) is normal in \(\operatorname{Gal}(E/F)\).
\end{exercise}

\subsection*{Looking Ahead}
\begin{exercise}[Looking Ahead] For \(n\mid m\), verify that the restriction
map
\[
\operatorname{Gal}(\mathbb F_{q^m}/\mathbb F_q)
\longrightarrow
\operatorname{Gal}(\mathbb F_{q^n}/\mathbb F_q)
\]
corresponds, under the Frobenius generators, to reduction
\[
\mathbb Z/m\mathbb Z\longrightarrow\mathbb Z/n\mathbb Z.
\]
\end{exercise}
\begin{exercise}[Looking Ahead] Explain why finite separable subextensions of
\(F^{\mathrm{sep}}/F\) correspond to open subgroups of \(G_F\).
\end{exercise}
